% from ega3-1-fr.tex
\oldpage[III]{104}
\begin{corollary}[2.3.2]
\label{III.2.3.2-fr}
Sous les hypothèses de (\hyperref[III.2.3.1-fr]{2.3.1}), soit
\[
\mathcal S_X=\bigoplus_{n\in\mathbb Z}\mathcal O_X(n),
\]
et soit $\mathcal F$ un $\mathcal S_X$-Module gradué quasi-cohérent de
type fini. Alors $\bbGamma_*(\mathcal F)$ vérifie la condition
\textbf{(TF)}.
\end{corollary}

\begin{proof}
On a vu dans la démonstration de
(\hyperref[III.2.3.1-fr]{2.3.1}) que $X$, qui est isomorphe à
$\operatorname{Proj}(\mathcal S^{(d)})$ (II, 3.1.8), est de type fini sur
$Y$ (II, 3.4.1). Il résulte alors de (II, 8.14.9) que $\mathcal F$ est
isomorphe à un $\mathcal S_X$-Module gradué de la forme
$\shProj(\mathcal M)$, où $\mathcal M$ est un $\mathcal S$-Module gradué
quasi-cohérent de type fini. En vertu de
(\hyperref[III.2.3.1-fr]{2.3.1}), $\bbGamma_*(\mathcal F)$ est
\textbf{(TN)}-isomorphe à $\mathcal M$ et par suite vérifie
\textbf{(TF)}.
\end{proof}

\begin{env}[2.3.3]
\label{III.2.3.3-fr}
\emph{Scholie.} --- Soient $Y$ un préschéma noethérien, $\mathcal S$ une
$\mathcal O_Y$-Algèbre graduée vérifiant les conditions de
(\hyperref[III.2.3.1-fr]{2.3.1}) et $X=\operatorname{Proj}(\mathcal S)$.
Soient $K_{\mathcal S}$ la catégorie abélienne des $\mathcal S$-Modules
gradués quasi-cohérents vérifiant \textbf{(TF)}, $K'_{\mathcal S}$ la
sous-catégorie de $K_{\mathcal S}$ formée des $\mathcal S$-Modules
vérifiant \textbf{(TN)}~; enfin, soit $K_X$ la catégorie des
$\mathcal S_X$-Modules gradués quasi-cohérents de type fini $\mathcal F$
(ce qui revient à dire, puisque $\mathcal S_X$ est périodique
(II, 8.14.4 et 8.14.12), que les $\mathcal F_i$ sont des
$\mathcal O_X$-Modules cohérents). Alors les foncteurs
\[
\mathcal M\longmapsto\shProj(\mathcal M)
\quad\text{dans }K_{\mathcal S},
\qquad
\mathcal F\longmapsto\bbGamma_*(\mathcal F)
\quad\text{dans }K_X
\]
définissent, en vertu de (II, 8.14.8 et 8.14.10) et de
(\hyperref[III.2.3.2-fr]{2.3.2}), une équivalence (T, I, 1.2) de la
catégorie quotient $K_{\mathcal S}/K'_{\mathcal S}$ (T, I, 1.11) avec la
catégorie $K_X$. Lorsque $\mathcal S$ est engendrée par $\mathcal S_1$,
on peut remplacer $K_X$ par la catégorie des $\mathcal O_X$-Modules
cohérents (II, 8.14.12).
\end{env}

\begin{proposition}[2.3.4]
\label{III.2.3.4-fr}
Soit $Y$ un préschéma noethérien.
\begin{enumerate}
\item[{\rm(i)}] Soit $\mathcal S$ une $\mathcal O_Y$-Algèbre graduée à
degrés positifs, quasi-cohérente et de type fini. Soient
$X=\operatorname{Proj}(\mathcal S)$, et
\[
\mathcal S_X=\shProj(\mathcal S)
=\bigoplus_{n\in\mathbb Z}\mathcal O_X(n).
\]
Alors $\mathcal S_X$ est une $\mathcal O_X$-Algèbre graduée périodique
(II, 8.14.12) dont les composants homogènes
$(\mathcal S_X)_n=\mathcal O_X(n)$ sont des $\mathcal O_X$-Modules
cohérents, et si $d>0$ est une période de $\mathcal S_X$,
$(\mathcal S_X)_d=\mathcal O_X(d)$ est un $\mathcal O_X$-Module
inversible $Y$-ample. En outre, l'homomorphisme canonique
\[
\alpha:\mathcal S\longrightarrow\bbGamma_*(\mathcal S_X)
\]
est un \textbf{(TN)}-isomorphisme.
\item[{\rm(ii)}] Inversement, soit $q:X\to Y$ un morphisme projectif, et
soit $\mathcal S'$ une $\mathcal O_X$-Algèbre graduée, dont les composants
homogènes $\mathcal S'_n$ $(n\in\mathbb Z)$ sont des
$\mathcal O_X$-Modules cohérents, et qui admet une période $d>0$ telle que
$\mathcal S'_d$ soit un $\mathcal O_X$-Module inversible ample pour $q$.
Alors
\[
\mathcal S=\bigoplus_{n\geq0}q_*(\mathcal S'_n)
\]
est une $\mathcal O_Y$-Algèbre graduée à degrés positifs quasi-cohérente
et de type fini, et il existe un $Y$-isomorphisme
\[
r:X\xrightarrow{\sim}\operatorname{Proj}(\mathcal S)
\]
tel que $r^*(\shProj(\mathcal S))$ soit isomorphe (en tant que
$\mathcal O_X$-Algèbre graduée) à $\mathcal S'$.
\end{enumerate}
\end{proposition}

\begin{proof}
(i) Toutes les assertions ont pratiquement déjà été démontrées, la
dernière n'étant autre qu'un cas particulier de
(\hyperref[III.2.3.2-fr]{2.3.2}). Le fait que $\mathcal S_X$ est
périodique a été vu en (II, 8.14.14) et le fait qu'il y a une période
$d>0$ telle que $\mathcal O_X(d)$ soit inversible et $Y$-ample n'est
autre que (II, 4.6.18). Enfin, pour $0\leq k<d$,
$(\mathcal S_X)^{(d,k)}$ est un $(\mathcal S_X)^{(d)}$-Module de type fini
(II, 8.14.14), donc chacun des $(\mathcal S_X)_n$ est un
$\mathcal O_X$-Module quasi-cohérent de type fini, en vertu de
(II, 2.1.6, (ii))~; la question étant locale, comme $\mathcal O_X$ est
cohérent, il en est de même des $(\mathcal S_X)_n$.

(ii) Quitte à remplacer la période $d$ par un de ses multiples, on peut
supposer que $\mathcal L=\mathcal S'_d$ est un $\mathcal O_X$-Module
\emph{très ample} relativement à $q$ (II, 4.6.11). On a en outre
\[
\mathcal S'^{(d)}=\bigoplus_{n\in\mathbb Z}\mathcal L^{\otimes n}
\quad\text{par hypothèse, donc}\quad
\mathcal S^{(d)}=\bigoplus_{n\geq0}q_*(\mathcal L^{\otimes n})~;
\]
on sait (II, 3.1.8 et 3.2.9) qu'il y a un $Y$-isomorphisme $s$ de
$X'=\operatorname{Proj}(\mathcal S)$ sur
$X''=\operatorname{Proj}(\mathcal S^{(d)})$ tel que
